\documentclass[11pt]{article}
\usepackage[a4paper,margin=27mm]{geometry}
\usepackage{amsmath,amssymb,mathtools}
\usepackage{fontspec,unicode-math}
\setmainfont{Latin Modern Roman}
\setmathfont{Latin Modern Math}
\usepackage{xurl}
\usepackage[hidelinks]{hyperref}
\newcommand{\sref}[2]{\hyperlink{source:#1:#2}{[S#2]}}
\newcommand{\cataloguescope}[1]{\paragraph{Mathematical question.}#1}
\hypersetup{pdftitle={Problem statement}}
\begin{document}
\section{Are there infinitely many Euclid primes}\label{statement}

\subsection{Definitions and mathematical statement}

\paragraph{Shared notation.}$\mathbb N=\{1,2,\ldots\}$, and $\mathbb Z,\mathbb Q,\mathbb R,\mathbb C$ denote the integers, rationals, reals and complex numbers. Any other conventions are specified locally.


Let $p_1=2<p_2=3<\cdots$ enumerate primes and $P_n=\prod_{j=1}^{n}p_j$ for $n\ge1$. Put $E_n=P_n+1$, the $n$th Euclid number.
\paragraph{Statement recorded in the supplied source.}
Is $E_n$ prime for infinitely many $n\ge1$? \sref{statement}{2}
\subsection{Short English statement}
Are infinitely many products of all primes up to a cutoff, plus one, themselves prime?
\subsection{Sources}
\begin{description}
\item[\hypertarget{source:statement:1}{S1}] ULAM.AI and the UnsolvedMath Contributors, UnsolvedMath: A Curated Collection of Open Mathematics Problems, catalogue entry. CC BY 4.0. \url{https://huggingface.co/datasets/ulamai/UnsolvedMath}.
\item[\hypertarget{source:statement:2}{S2}] Wikipedia contributors, \emph{List of unsolved problems in mathematics}, revision 1366636767 (2026) Number theory, Are there infinitely many Euclid primes. \url{https://en.wikipedia.org/w/index.php?title=List_of_unsolved_problems_in_mathematics&oldid=1366636767#Number_theory}.
\item[\hypertarget{source:statement:3}{S3}] OEIS Foundation and contributors, \emph{The On-Line Encyclopedia of Integer Sequences}, A006862: Euclid numbers: 1 + product of the first n primes., official data snapshot \texttt{693b5499c9d2} (2026) Name, definition and comments. \url{https://raw.githubusercontent.com/oeis/oeisdata/693b5499c9d248f570dd2d979e8bd1d7eecf689b/seq/A006/A006862.seq}.
\item[\hypertarget{source:statement:4}{S4}] OEIS Foundation and contributors, \emph{The On-Line Encyclopedia of Integer Sequences}, A014545: Numbers k such that the k-th Euclid number A006862(k) = 1 + (Product of first k primes) is prime., official data snapshot \texttt{693b5499c9d2} (2026) Name, definition and comments. \url{https://raw.githubusercontent.com/oeis/oeisdata/693b5499c9d248f570dd2d979e8bd1d7eecf689b/seq/A014/A014545.seq}.
\item[\hypertarget{source:statement:5}{S5}] Wikipedia contributors, \emph{Euclid number}, revision 1321323502 (2026) Definition and Unsolved problems. \url{https://en.wikipedia.org/w/index.php?title=Euclid+number&oldid=1321323502}.
\end{description}
\label{end:statement}
\end{document}
